\section{결론}
\label{sec:결론}

본 논문에서는 급격한 토폴로지 변동과 한정된 컴퓨팅 및 통신 자원을 지닌 Flying Ad Hoc Networks (FANETs) 환경을 극복하기 위해 글로벌-로컬 정책 증류 및 예측적 위험 스위칭 기반의 분산형 라우팅 프레임워크인 \method를 제안하였다. 제안 기법은 중앙 집중식 오프라인 학습 단계에서 전역 토폴로지 그래프 정보를 활용하는 privileged 교사 강화학습 모형을 구축한 뒤, 그 행동 양식을 1-hop 로컬 관측 정보만을 사용하는 경량 학생 모델에 증류함으로써 배포 단계에서의 전역 정보 공유 비용을 제거하였다. 

또한, 런타임 실행 단계에서 가벼운 지리적 포워딩을 수행하면서 링크 수명, onward 안정성, 버퍼 대기열 등의 지표를 모니터링하여 위험 임계 상황에서만 예측 지향 경로 선택기(PRS)를 선택 구동시키는 하이브리드 제어 루프를 제안하였다. 시뮬레이션 평가 결과, 제안 기법은 GPSR, Predictive Geo, Evo-QGeo, IQMR, DRAMA와 같은 대표적인 외부 라우팅 알고리즘군과 비교하여 패킷 전달율(PDR)과 데드라인 만족 전송률을 유의미하게 향상시켰으며, 로컬 구동에 필요한 입력 풋프린트 오버헤드와 통신 제어 오버헤드를 낮게 조율함을 입증하였다. 향후 연구로는 ns-3 시뮬레이터를 활용한 대규모 교차 검증 및 다중 무인기 실물 테스트베드를 통한 런타임 실효성 검증을 진행할 계획이다.
